\documentclass[11pt]{article}
\usepackage[margin=1.1in]{geometry}
\usepackage{amsmath,amssymb,amsthm}
\usepackage[T1]{fontenc}
\usepackage[expansion=false]{microtype}
\usepackage{booktabs}
\usepackage{hyperref}

\newtheorem*{prop}{Proposition}

\title{The Order of Arrival:\\[4pt]
\large Density, Balance, and Transient Structure in the Evolution of Random Graphs}
\author{Flyxion\\ \small Independent Researcher}
\date{September 2026}

\begin{document}
\maketitle

\begin{abstract}
Erd\H{o}s and R\'enyi's 1960 paper on the evolution of random graphs shows that as edges are added uniformly at random to $n$ vertices, small structures arrive in a fixed order, each at a definite scale of the edge count relative to $n$. This essay reads those results as a study of what an edge budget permits. Three points organize the reading. The timing of a structure's arrival is set by its density, and specifically by the densest part of it. The exponent of arrival and the abundance on arrival are governed by different features: density fixes the first, combinatorial multiplicity fixes the second. And the properties the process exhibits divide into those that can only be gained and those that can be gained and then lost. Containing a tree is of the first kind; having a tree stand alone as a component is of the second. The second kind has no threshold, only a window, and the window marks the difference between a structure existing and a structure being individuated.
\end{abstract}

\section{Introduction}

Begin with $n$ labeled vertices and no edges. Choose a pair not yet joined, uniformly at random, join it, and repeat. After $N$ steps the result is the random graph $G_{n,N}$. Nothing in the rule refers to trees, cycles, or cliques, and no step aims at any structure. Yet as $N$ grows, structures arrive in a fixed order, each at a scale of $N$ that can be computed in advance, and each with a limiting distribution that can be written in closed form.

The order is not chosen by anything. It falls out of counting: how many ways a structure can be placed among $n$ vertices, set against how many edges it consumes. That trade is the entire mechanism, and the 1960 results are its consequences worked out with care. What makes the paper worth rereading is not only the answers but the shape of the reasoning. It is an account of structure in which nothing is built, only permitted, and in which permission is a function of a single ratio.

\section{The Process and Its Clock}

For a fixed graph $H$ with $k$ vertices and $l$ edges, the natural question is when copies of $H$ begin to appear. A \emph{threshold function} $A(n)$ answers it in the asymptotic sense: if $N(n)/A(n)\to 0$ then $G_{n,N}$ almost surely contains no copy of $H$, and if $N(n)/A(n)\to\infty$ it almost surely contains one.

The heuristic behind every threshold in the paper is a first-moment count. Write $p \approx 2N/n^2$ for the chance that a given pair is joined. The number of placements of $H$ on labeled vertices is of order $n^k$, and each placement is realized with probability about $p^l$. The expected number of copies is therefore of order
\[
n^k p^l \;\asymp\; n^k \left(\frac{N}{n^2}\right)^{l} \;=\; N^l\, n^{k-2l}.
\]
This is of order one exactly when
\[
N \asymp n^{\,2-k/l}.
\]
The exponent depends on $H$ only through the ratio $l/k$, the number of edges per vertex. Everything else about the shape of $H$ drops out of the exponent. The relevant clock is not $N$ but $N$ measured in units of $n^{2-k/l}$, and different structures run on different clocks.

A convenient restatement: with $d = l/k$, the threshold exponent is $2 - 1/d$. Structures with $d<1$ arrive before $N \asymp n$, structures with $d=1$ arrive at $N\asymp n$, and structures with $d>1$ arrive after it, with the exponent climbing toward $2$ as density grows. The whole map of arrival is a monotone function of one number.

\section{Balance and the Densest Part}

The first-moment argument has a gap. The expected number of copies of $H$ can grow without bound while $H$ is almost surely absent. This happens when $H$ contains a subgraph $H'$ denser than itself. Since $H'$ has a larger edge-to-vertex ratio, it has a later threshold, and $H$ cannot appear before its own subgraph does. In the intervening range the expectation is carried by rare graphs that happen to contain $H'$ and therefore contain many copies of $H$ built around it.

A concrete case: attach one pendant edge to $K_4$. The whole graph has $5$ vertices and $7$ edges, so the naive exponent is $2-5/7 = 9/7$. But $K_4$ alone has ratio $6/4$, exponent $2 - 4/6 = 4/3$, and $4/3 > 9/7$. For $N$ between $n^{9/7}$ and $n^{4/3}$ the expected count of the whole configuration diverges while the configuration is almost surely absent.

Erd\H{o}s and R\'enyi avoid the gap by restricting to \emph{balanced} graphs: graphs in which no subgraph has a higher average degree than the whole. For a connected balanced graph with $k$ vertices and $l\ge k-1$ edges they prove that $n^{2-k/l}$ is the threshold. The restriction is not a technicality to be waved away. It encodes the fact that arrival waits on the tightest local constraint. Bollob\'as later removed the restriction by replacing $l/k$ with the maximum of $l_{H'}/k_{H'}$ over all subgraphs $H'$, which makes the principle explicit: a structure arrives when its densest part can.

This is the first organizing point. The edge budget acts as a constraint, and a structure is permitted when every part of it is permitted. The binding part is the densest one, and a structure with a dense core and sparse periphery is timed entirely by the core.

\section{The Ladder of Trees}

A tree on $k$ vertices has $l = k-1$ edges and is balanced, so its threshold is
\[
N \asymp n^{(k-2)/(k-1)}.
\]
The exponents run $0, \tfrac12, \tfrac23, \tfrac34, \dots$, increasing toward $1$ without reaching it. Single edges appear as soon as $N\to\infty$; paths of length two at $\sqrt{n}$; trees on four vertices at $n^{2/3}$. Every tree order has its rung, and all infinitely many rungs are compressed into the range below $N \asymp n$.

Because only $k$ and $l$ enter the exponent, all trees of the same order arrive together. A path and a star on the same number of vertices share a threshold. Shape is invisible at the level of timing.

Shape reappears in the constant. When $N \sim \rho\, n^{(k-2)/(k-1)}$ with $\rho>0$ fixed, the number of tree components of order $k$ converges to a Poisson distribution with mean
\[
\lambda = \frac{(2\rho)^{k-1}\,k^{k-2}}{k!}.
\]
The factor $k^{k-2}$ is Cayley's count of labeled trees on $k$ vertices. It enters as multiplicity: the number of distinct ways the order-$k$ tree condition can be met. The factor $(2\rho)^{k-1}$ carries the edge cost, and $k!$ removes the labeling of the vertex set. Directly from the first-moment count, with $p = 2\rho\, n^{(k-2)/(k-1)-2}$,
\[
\binom{n}{k} k^{k-2} p^{k-1} \;\to\; \frac{k^{k-2}}{k!}(2\rho)^{k-1},
\]
since $p^{k-1} = (2\rho)^{k-1} n^{-k}$ exactly cancels $n^k$.

This separation is the second organizing point. The exponent answers \emph{when}, and it depends on density alone. The constant answers \emph{how many}, and it depends on how many configurations share that density. Two structures can arrive on the same clock and in very different quantities.

\section{Cycles Arrive Together}

A cycle on $k$ vertices has $l=k$ edges, so its exponent is $2 - k/k = 1$ for every $k\ge 3$. The ladder of trees ends, and every cycle length shares a single threshold, $N\asymp n$. Triangles, pentagons, and cycles of length one hundred all begin to appear on the same scale.

Length is again visible only in the constant. For $N\sim cn$, the number of $k$-cycles converges to a Poisson distribution with mean
\[
\lambda_k = \frac{(2c)^k}{2k},
\]
where $2k$ counts the rotations and reflections that fix a labeled cycle. For $c<1/2$ we have $2c<1$, so the means decay geometrically in $k$ and short cycles dominate. Summing over lengths,
\[
\sum_{k\ge 3} \frac{(2c)^k}{2k} \;=\; -\tfrac12\log(1-2c) - c - c^2,
\]
which is finite for $c<1/2$ and diverges as $c\to 1/2$. This divergence is the arithmetic signature of the transition at $N\sim n/2$, where the giant component emerges. Below it, cycles of each length are rare and their total is controlled. At it, cycles of every length stop being rare at once.

The paper states the corresponding exact result: for $N\sim cn$ with $c<1/2$, the probability that $G_{n,N}$ contains no cycle at all tends to
\[
\sqrt{1-2c}\;e^{\,c+c^2}.
\]
This is precisely $\exp\!\big(-\sum_{k\ge3}\lambda_k\big)$, the value the independent Poisson picture predicts, with $e^{c+c^2}$ removing the phantom terms for $k=1$ and $k=2$ that the logarithm would otherwise include. As $c\to 1/2$ the expression tends to $0$, and a cycle becomes almost certain.

Complete graphs continue the map above $n$. $K_k$ has ratio $(k-1)/2$ and exponent $2 - 2/(k-1)$: the triangle at $n$, $K_4$ at $n^{4/3}$, $K_5$ at $n^{3/2}$, approaching $n^2$. The overall picture is a single line of exponents from $0$ to $2$, with trees filling the interval below $1$, cycles sitting exactly at $1$, and denser graphs spread above it.

\begin{center}
\begin{tabular}{@{}lccc@{}}
\toprule
Structure & $k$ & $l$ & Threshold \\
\midrule
Tree of order $k$ & $k$ & $k-1$ & $n^{(k-2)/(k-1)}$ \\
Cycle of order $k$ & $k$ & $k$ & $n$ \\
Complete graph $K_k$ & $k$ & $\binom{k}{2}$ & $n^{2-2/(k-1)}$ \\
Isolated cycle of order $k$ & $k$ & $k$ & none (window) \\
\bottomrule
\end{tabular}
\end{center}

\section{Windows}

Some structures in the paper are not only gained. Consider an \emph{isolated} $k$-cycle, one whose vertices carry no other edges. Its expected count picks up a second factor, the probability that none of the roughly $kn$ other pairs touching the cycle are joined, which is about $e^{-2ck}$. For $N\sim cn$ the count is asymptotically Poisson with mean
\[
\mu_k = \frac{(2c\,e^{-2c})^k}{2k}.
\]
The function $c\mapsto 2c\,e^{-2c}$ rises, peaks at $c = 1/2$ with value $e^{-1}$, and falls. So the probability of seeing an isolated $k$-cycle rises, is greatest exactly at the critical point, and then declines to zero as the cycles are absorbed into the growing giant component.

Two constraints pull in opposite directions here. Being a cycle requires edges to be present. Being isolated requires edges to be absent. The first is favored by a larger budget and the second by a smaller one, and the structure is visible only in the range where both can hold.

Trees behave the same way at a later scale. Isolated trees of order $k$ persist long after their arrival and begin to disappear around $N\sim \frac{1}{2k}\,n\log n$. At the finer scale
\[
N \sim \tfrac{1}{2k}\,n\log n + \tfrac{k-1}{2k}\,n\log\log n + yn,
\]
their number is again asymptotically Poisson, now with mean $e^{-2ky}/(k\cdot k!)$. Because the leading term $\frac{1}{2k}n\log n$ shrinks as $k$ grows, large isolated trees are absorbed first and small ones last. The last to go are isolated vertices, at $N\sim\frac12 n\log n$, which is the threshold for connectivity. The process has an order of disappearance as well as an order of arrival, and it runs from large pieces to small.

This is the third organizing point. Containing a fixed graph is a \emph{monotone} property: adding an edge never destroys it. Bollob\'as and Thomason proved that every monotone property of this kind has a threshold. Having an isolated copy of a fixed graph is not monotone, since an added edge can destroy it, and such properties need not have a threshold at all. What they have instead is a window. The statement that isolated cycles have no ordinary threshold is not a curiosity of cycles; it is what non-monotonicity looks like.

\section{Containment and Component}

A tree that has been absorbed into the giant component has not gone anywhere. Every edge it had is still present, and it is still a subgraph. What has been lost is its separability: the fact that it could be read off as a component without cutting anything.

The process only adds. Nothing is ever deleted. Yet components are lost continually, by merger, and the number of components falls monotonically from $n$ to $1$. This is loss by addition. The identity of a component depends on a boundary condition, that no edge leaves it, and a boundary condition is exactly the kind of thing that additions violate. The structure survives; the constraint that made it an individual does not.

The distinction generalizes beyond graphs. Whether something exists and whether it can be picked out as a separate object are different questions, answered by different kinds of condition. The first is typically a presence condition and tends to be stable under accumulation. The second is typically a closure condition and tends to be eroded by accumulation. A process that only accumulates will therefore show a characteristic pattern: structures become more numerous and more permanent as parts, while becoming less available as wholes.

\section{Four Transfers}

The random-graph process is tempting to export because its pictures are so portable. A sparse field becomes connected; local pieces join a macroscopic component; an apparently continuous increase in edges produces a qualitative change in organization. Physical computers, biological tissues, and social cascades all admit pictures of this kind. But visual resemblance is not yet a mathematical correspondence. The transition at $N\sim n/2$ belongs to a specific ensemble: unweighted edges are sampled uniformly, vertices have no geometry, and the graph grows without deletion or preference. A collagen network, an Ising machine, and a repost cascade violate nearly every one of those assumptions.

What transfers more reliably is therefore not the critical number but the separation of questions that produces it. Four questions recur. What resource makes a structure admissible? Which local part imposes the binding requirement? How many realizations become available at the same cost? And does the structure remain separately identifiable after further connections are added? In the random graph these are answered by density, maximum subgraph density, combinatorial multiplicity, and closure. Elsewhere the quantities differ, but the separation remains useful.

\subsection{Physical Ising Machines: the Cost of the Dense Part}

An Ising formulation assigns binary variables $s_i\in\{-1,+1\}$ and couplings $J_{ij}$, with an objective of the form
\[
E(s)=-\sum_{i<j}J_{ij}s_is_j-\sum_i h_i s_i.
\]
The variables scale linearly with the problem size, but the number of possible pairwise couplings scales quadratically. A statement that a machine contains $n$ spins consequently says little about which $n$-variable Hamiltonians it can realize directly. The operative resource is not spin count alone but the topology, programmability, precision, and physical fan-out of the couplings.

This gives a concrete analogue of balance. A logical problem may be sparse on average while containing a subproblem whose variables must be coupled densely. That dense core controls the embedding. Sparse variables attached around it do not lower the core's routing demand, just as a pendant edge attached to $K_4$ does not advance the arrival of the resulting graph. The whole construction must wait until its most demanding part can be represented. Hardware claims based only on the total number of variables can therefore conceal the relevant constraint in much the same way that average density can conceal an unbalanced subgraph.

Embedding changes the accounting rather than abolishing it. A dense logical graph can be represented on a sparse physical graph by chains, auxiliary variables, time multiplexing, or repeated measurement, but each device converts missing connectivity into another expenditure. The natural comparison between machines is thus not a single nominal scale. It is a map from logical topology to physical cost: spins consumed per logical variable, routing depth, precision lost, samples required, and the probability that the representation preserves the intended ground state. The dense part of the logical instance is where these costs are likely to bind.

Noise supplies a second, more limited comparison. Thermal, optical, or deliberately injected fluctuations can help a solver move among configurations, and annealing schedules can cross changes in the organization of an energy landscape. But this is not the Erd\H{o}s--R\'enyi giant-component transition. The control parameter, ensemble, and order parameter are different, and noise does not guarantee arrival at a ground state. The transferable lesson is methodological: before calling two abrupt changes the same phase transition, one must specify what is varying, what macroscopic quantity changes, and under which distribution the claimed threshold is typical.

\subsection{Mechanical Networks: Connection and Individuation}

The extracellular matrix is also a network, but not an Erd\H{o}s--R\'enyi graph. Its links have length, direction, stiffness, and history. Fibres align under load; crosslinks are chemically heterogeneous; deposition and degradation occur locally; cells both respond to and alter the network through which forces reach them. Mechanical connectivity is consequently spatial, weighted, correlated, and active.

Even so, the distinction between containing a structure and sustaining it as a component carries over with unusual force. A collagen arrangement can remain materially present while losing the mechanical boundary that once allowed it to move, deform, or remodel as a distinguishable unit. Additional crosslinks may increase transmission across the tissue while decreasing local independence. Nothing need be removed. The earlier configuration survives inside a more connected matrix, but it no longer exists in the same sense as a mechanically individuated object. This is the biological form of loss by addition.

The movement--stress--deformation--remodelling loop can therefore be described without pretending that fibrils are sampled random edges. Movement supplies patterned load; load selects which connections bear stress; deformation changes local geometry; remodelling stabilizes or releases some of those relations. The process alters both the paths along which force can percolate and the boundaries within which deformation can remain local. A tissue may become more globally coherent precisely by becoming less locally permissive.

This interpretation sharpens the history--plasticity axis $\Lambda$. An increase in accumulated history $H$ and a decrease in remodelling plasticity $P$ may produce a qualitative mechanical transition, but the name \emph{phase transition} should be earned rather than supplied by analogy. It requires a stated control parameter, a measurable order parameter, and evidence that the change is not merely gradual stiffening. Candidate observables might include system-spanning force transmission, rigidity, relaxation time, or the fraction of perturbations that remain locally recoverable. The random-graph analogy then serves as a question generator: which connections create global transmission, which dense local motifs bind first, and which former units cease to be separable as connectivity accumulates?

The distinction also prevents a common inversion. Greater connectivity is not simply greater function, nor is lower connectivity simply greater plasticity. The effect depends on the topology and on the scale at which function is measured. A newly completed force path may permit coordinated tissue action while simultaneously closing a local degree of freedom. Global acquisition and local loss can be the same event viewed across two boundaries.

\subsection{Information Cascades: a Different Critical Process}

A repost cascade resembles a growing tree, but its growth rule differs fundamentally from random edge addition. The social network largely pre-exists the information unit; exposure is nonuniform; high-degree accounts have disproportionate reach; platform ranking intervenes; and every transmission occurs in time. The natural elementary model is therefore often a branching or epidemic process on a heterogeneous network rather than an Erd\H{o}s--R\'enyi graph being assembled from nothing.

For that reason, the appearance of a large cascade should not be identified with the $c=1/2$ giant-component threshold. In the convention used here, $N\sim cn$ gives mean degree approximately $2c$, and the giant component appears when that mean reaches one. A diffusion model has its own critical condition, usually expressed through an effective reproduction quantity or through the leading structure of a transmission operator. The threshold depends on exposure, adoption, network heterogeneity, and intervention. The number $1/2$ does not transfer.

What does transfer is the problem of inference before macroscopic individuation. During an early observation window $\Delta_o$, a monitoring system sees only a partial cascade: a certain breadth, depth, branching profile, velocity, and arrangement of participating accounts. It must decide whether these observations belong to a process likely to terminate as a small tree or to one capable of sustaining growth into a high-amplification tail. An index such as DDVRI is useful to the extent that it estimates that latent growth regime rather than merely redescribing current size.

Here the separation between timing and multiplicity becomes important. Many shallow realizations may share the same early size while having very different continuation probabilities. Conversely, two early cascades with different shapes may be equivalent with respect to a properly specified diffusion model. The predictive question is not simply how much structure has already appeared, but which observed feature is binding for further growth. An apparent dense core of mutually reinforcing accounts, a bridge into a new audience, or repeated exposure through several independent paths may matter more than the average branching count.

Closure matters here as well, though in an altered form. A cascade is often treated as an individual information event because its posts can be assigned to one source claim or lineage. As quotations, screenshots, corrections, and mutations connect it to neighbouring narratives, that attributional boundary can disappear even while every original post remains. Accumulation increases the observable record but erodes the separability of the object being measured. A monitoring system must therefore state whether it predicts the continuation of a message, a claim, a semantic family, or a socially identified event. Otherwise its target can dissolve while its count continues to rise.

\subsection{Empirical Evaluation: Expectation Is Not Admissibility}

The most exact transfer is methodological. In the unbalanced-graph example, the expected number of copies of a structure can diverge while the probability of seeing even one copy tends to zero. The expectation is dominated by rare samples in which the dense prerequisite occurs and supports many extensions. A summary statistic reports abundance without reporting whether the condition that makes ordinary realization possible has been met.

Benchmark evaluation can fail in the same shape. Suppose an optimization system is tested on random number-partitioning instances drawn from a narrow integer range. Such a distribution may contain many exact or near-exact partitions, making successful outputs combinatorially abundant. A large success rate then demonstrates performance on that ensemble, but it does not by itself establish reliable solution of distributions in which satisfying assignments are sparse, nor does it certify that any reported solution is globally optimal. The number of opportunities to look correct and the difficulty of satisfying the binding constraint have been confounded.

This is not an argument against probabilistic evidence. Erd\H{o}s and R\'enyi's results are themselves precise probabilistic statements. It is an argument for naming the quantifier. A probable threshold says that a property holds with probability tending to one under a specified ensemble. A compulsory bound states what every graph satisfying stated conditions must contain. An average-case benchmark, a worst-case complexity claim, and an independently verified optimum answer different questions. None can silently substitute for another.

The same discipline separates verification from self-agreement. Repeated runs that converge on the same answer estimate stability under the chosen procedure. They do not establish correctness unless the procedure is tied to an independent witness: an exact solver, a checkable certificate, a valid lower bound meeting the candidate upper bound, or another criterion that cannot be satisfied merely because the system repeats itself. Multiplicity may make an output common; only an external condition makes the intended claim admissible.

The densest-part principle gives a further diagnostic. A benchmark suite may appear varied while all its instances avoid the local configurations that impose the real computational burden. Mean size, mean density, or total instance count can then conceal the absence of the binding substructure. Evaluation should identify the feature believed to control hardness and vary it directly. Otherwise the suite measures the abundance of easy realizations rather than the arrival of the claimed capability.

\section{What Transfers}

These four domains do not share one universal phase transition. They share a recurrent error: treating a global summary as though it disclosed the local condition that produced it. Spin count conceals coupling topology. Matrix density conceals anisotropy and mechanical boundaries. Cascade size conceals the process governing continuation. Benchmark accuracy conceals both instance distribution and the independence of verification.

The random-graph analysis supplies a compact corrective. First identify the budget or control parameter. Then find the local structure that imposes the latest or most expensive requirement. Keep the threshold of appearance separate from the number of realizations available after appearance. Finally, ask whether additional connections preserve the object of interest or merely preserve its parts while destroying the closure by which it was recognized.

This sequence is stronger than a general appeal to emergence. Emergence says that a macroscopic pattern arose from local relations. The order-of-arrival analysis asks why this pattern could arise at this scale, why another had to wait, why some patterns were plentiful when they arrived, and why further accumulation caused an apparently permanent object to disappear as a whole. It converts resemblance into a set of discriminations.

\section{What the Process Permits}

Read in this way, the evolution of random graphs is a one-parameter family of constraints. At each value of $N$ the edge budget determines which structures are expected, in what quantities, and in what configurations. Nothing is selected and nothing is assembled toward an end. The order of arrival follows from comparing multiplicity with cost, and the order of disappearance follows from comparing presence conditions with closure conditions.

Two cautions keep the reading honest. First, thresholds are statements about sequences of distributions as $n$ grows, not about any single graph. The ``evolution'' is a way of traversing a family of probability spaces, and the narrative of arrival and absorption is a description of how typical members of that family change with $N$. Second, the Erd\H{o}s--R\'enyi model has no geometry, no locality, and no preferential attachment. Real networks violate its assumptions in ways that change the numbers. The lesson carried over is not the numbers but the logic: timing is set by a binding cost, abundance by multiplicity, and individuation by a closure condition that accumulation may wear away. Density gives the exact cost only inside the random-graph model; elsewhere the task is to discover what plays its role.

The paper's quiet achievement is to show that a process with no preferences produces a strict ordering of structures, and that the ordering can be read from a single ratio. What arrives first is what costs least per vertex. What survives longest as a separate thing is what is smallest. Between those two facts lies most of the story.

\appendix
\section{The First-Moment Count and the Cycle Sum}

For $H$ with $k$ vertices, $l$ edges, and automorphism group of order $a(H)$, the expected number of copies of $H$ in $G(n,p)$ is
\[
\frac{(n)_k}{a(H)}\,p^l,
\]
where $(n)_k = n(n-1)\cdots(n-k+1)$. With $p \approx 2N/n^2$ and $k$ fixed this is asymptotic to $n^k (2N)^l n^{-2l}/a(H)$, which is bounded away from $0$ and $\infty$ exactly when $N\asymp n^{2-k/l}$. For a cycle $a = 2k$; for a labeled count of trees one sums over the $k^{k-2}$ labeled trees on a fixed $k$-set, which replaces $1/a(H)$ by $k^{k-2}/k!$.

For the cycle sum, $-\log(1-x) = \sum_{k\ge1} x^k/k$ gives
\[
\sum_{k\ge 3}\frac{(2c)^k}{2k} = \tfrac12\Big(-\log(1-2c) - 2c - 2c^2\Big),
\]
so that
\[
\exp\!\Big(-\sum_{k\ge3}\lambda_k\Big) = \sqrt{1-2c}\;e^{\,c+c^2},
\]
matching the limiting probability that $G_{n,N}$ is acyclic when $N\sim cn$, $c<1/2$.

\section*{References}
\small
\begin{itemize}\itemsep2pt
\item P.~Erd\H{o}s and A.~R\'enyi, On the evolution of random graphs, \emph{Publ.\ Math.\ Inst.\ Hung.\ Acad.\ Sci.} 5 (1960), 17--61.
\item B.~Bollob\'as, Threshold functions for small subgraphs, \emph{Math.\ Proc.\ Cambridge Philos.\ Soc.} 90 (1981), 197--206.
\item B.~Bollob\'as and A.~Thomason, Threshold functions, \emph{Combinatorica} 7 (1987), 35--38.
\item A.~Cayley, A theorem on trees, \emph{Quart.\ J.\ Pure Appl.\ Math.} 23 (1889), 376--378.
\end{itemize}

\end{document}
